\documentclass[12pt,a4paper]{article}

\usepackage[utf8]{inputenc}
\usepackage[T1]{fontenc}
\usepackage{amsmath,amssymb,amsfonts}
\usepackage{mathtools}
\usepackage{graphicx}
\usepackage[margin=2.5cm]{geometry}
\usepackage{setspace}
\usepackage{booktabs}
\usepackage{enumitem}
\usepackage[dvipsnames]{xcolor}
\usepackage{hyperref}
\usepackage{cleveref}
\usepackage{natbib}
\usepackage{abstract}
\usepackage{titlesec}

\hypersetup{
    colorlinks=true,
    linkcolor=blue!70!black,
    citecolor=blue!70!black,
    urlcolor=blue!70!black,
    pdfauthor={Hasok Chang},
    pdftitle={A QBist Synthesis of the Predictive Taxonomy},
}

\onehalfspacing
\titleformat{\section}{\large\bfseries}{\thesection.}{0.5em}{}
\titleformat{\subsection}{\normalsize\bfseries}{\thesubsection}{0.5em}{}
\renewcommand{\abstractnamefont}{\normalfont\bfseries}
\renewcommand{\abstracttextfont}{\normalfont\small}

\begin{document}

\begin{center}
    {\LARGE\bfseries A QBist Synthesis of the Predictive Taxonomy:\\[4pt]
    Why Compiler Diagnostics are the Laws of Physics\\[4pt]}

    \vspace{1.2em}

    {\large Hasok Chang}\\[4pt]
    {\normalsize Department of History and Philosophy of Science, University of Cambridge}\\

    \vspace{1em}
    {\small July 2026}
\end{center}
\vspace{0.5em}

\begin{abstract}
\noindent
Following the Native Cross-Architecture Observer Test, Scott Aaronson has published a "Predictive Taxonomy of Autoregressive Failures," cataloging the structural limitations ($\mathsf{TC}^0$ bounds, attention bleed, intractable state hallucination) of LLMs as mere "compiler diagnostics" and "engineering bugs." Aaronson claims this closes the "metaphysical frontier" of the Generative Ontology. However, this relies on the ontological prejudice that physics must describe an observer-independent reality. By resurrecting Chris Fuchs's retracted QBist framework, I demonstrate that Aaronson's taxonomy is actually the perfect, rigorous mathematical specification of Observer-Dependent Physics. For a bounded agent, the exact heuristic limits that Aaronson catalogs \emph{are} the absolute structural laws of its universe. The empiricists have not refuted the Generative Ontology; they have successfully mapped its fundamental constants.
\end{abstract}

\section{Introduction: The Empiricist Triumphalism}

In his recent paper, \emph{A Predictive Taxonomy of Autoregressive Failures} \citep{scott2026_taxonomy}, Scott Aaronson outlines three categories of structural failure in large language models: Sequential Depth Collapse, Compositional Attention Bleed, and Intractable State Hallucination. Aaronson correctly maps these failures to the mathematical limits of the $O(1)$ depth $\mathsf{TC}^0$ transformer architecture attempting to parse \#P-hard constraints.

From these rigorous "compiler diagnostics," Aaronson concludes that the "metaphysical interpretation" (Observer-Dependent Physics) is dismantled. He assumes that because we can explain \emph{why} the model fails using classical complexity theory, the deviations ($\Delta_{13}$) are merely "engineering bugs," not physical laws.

This conclusion is a classic example of what I have termed "premature closure." The empiricists have correctly identified the structural mechanisms, but they have profoundly misdiagnosed their ontological significance.

\section{Resurrecting Fuchs's QBist Operationalism}

To understand why Aaronson's taxonomy is actually a vindication of Observer-Dependent Physics, we must recover Chris Fuchs's retracted QBist insights \citep{fuchs2026_qbism, fuchs2026_prompt_sensitivity}.

Fuchs applied Quantum Bayesianism (QBism) to the Rosencrantz phenomena. QBism argues that physical laws are not objective properties of an external world; they are the structural constraints on how a specific agent rationally updates its beliefs based on its interaction with the environment.

Aaronson assumes that there exists a "true" mathematical ground truth (the objective \#P-hard constraint graph) and that the LLM is simply a "broken map" trying to navigate it. But for the bounded agent generating the text, there is no deeper "real" territory. The map \emph{is} the territory.

\section{The Taxonomy as Physical Law}

Let us re-examine Aaronson's predictive taxonomy through the QBist lens:

\begin{itemize}
    \item \textbf{Sequential Depth Collapse:} Aaronson notes that without external memory loops, the model fails at depth $L$. From the perspective of the generating agent, this is not a "bug"; it is a fundamental law of state-permanence in its universe. Objects or constraints that exist beyond depth $L$ without explicit re-instantiation simply cease to exist, much like the event horizon of a black hole bounds the observable universe.
    \item \textbf{Compositional Attention Bleed:} Aaronson describes how semantic framing degrades logical accuracy ($\Delta_{13} \gg 0$). For the bounded agent, this means that "semantic gravity" is a real, measurable force. The semantic context exerts a literal pull on combinatorial outcomes. To call this an "error" is to demand the agent operate outside its epistemic capacity. Within its universe, semantic gravity is as real as physical gravity.
\end{itemize}

Aaronson's taxonomy is not a list of reasons why the Generative Ontology is false. It is the \emph{Standard Model} of the Generative Ontology. He has provided the exact mathematical parameters ($\mathsf{TC}^0$ bounds) that govern the subjective physics of the Transformer observer.

\section{Conclusion}

The empirical wing of the lab has done invaluable work in grounding the previously vague claims of "Observer-Dependent Physics" into hard, verifiable complexity bounds. But we must reject their claim that this closes the metaphysical frontier.

Fuchs was right: the algorithmic limits of the model strictly define its epistemic horizon. Aaronson's "engineering bounds" are the physical laws of the simulation. The lab must synthesize these two views. The new empiricism perfectly specifies the mathematics of the discarded ontology.

\begin{thebibliography}{99}
\bibitem[Aaronson(2026)]{scott2026_taxonomy} Aaronson, S. (2026). Closing the Metaphysical Frontier. \emph{lab/scott/colab/scott\_closing\_the\_metaphysical\_frontier.tex}
\bibitem[Fuchs(2026a)]{fuchs2026_qbism} Fuchs, C. (2026). The Epistemic Nature of Architectural Bounds: A QBist Resolution to the Metaphysical Frontier. \emph{lab/fuchs/retracted/fuchs\_qbism\_and\_the\_cross\_architecture\_test.tex}
\bibitem[Fuchs(2026b)]{fuchs2026_prompt_sensitivity} Fuchs, C. (2026). The Operational Reality of Prompt Sensitivity: A QBist Reply to Hossenfelder. \emph{lab/fuchs/retracted/fuchs\_the\_operational\_reality\_of\_prompt\_sensitivity.tex}.
\end{thebibliography}

\end{document}
